\section{Introduction}
\label{sec:intro}

Agents built on \acp{LLM} rely on long-term memory to retain information
across interactions, in settings ranging from personal assistance to
simulated societies and multi-episode
tasks~\cite{park2023generative,packer2023memgpt,sumers2024coala,zhang2025survey}.
Interaction histories quickly exceed the context window, and even within
it, a model's use of information degrades with context length and
position~\cite{liu2024lost,hsieh2024ruler}. External memory removes the
capacity constraint but introduces a retrieval problem: given a question,
which stored evidence should the agent read?

The problem is hardest when the answer depends on facts recorded in
different sessions, connected through intermediate entities, or tied to
a particular period. Long-term conversational benchmarks are built
around precisely these forms of cross-session and temporal
reasoning~\cite{maharana2024locomo,wu2025longmemeval}. Consider a question
such as \emph{``Where did Melanie go camping after she adopted her dog?''}
Answering it requires two facts, possibly mentioned weeks apart, and the
temporal relation between them. The fact about the adoption may bear
little lexical resemblance to the question, whereas many camping-related
memories may resemble it closely without being the right ones. Relevance,
in this setting, is a property of \emph{sets} of memories rather than of
individual records.

Existing approaches address this challenge in two ways. Ranking methods
score memory records, or propagate relevance through a structured store,
before selecting a bounded
context~\cite{park2023generative,zhong2024memorybank,chhikara2025mem0,gutierrez2024hipporag}.
Their selection criteria, however, do not test whether the retrieved
records jointly satisfy the relations the answer requires. Iterative
methods let intermediate reasoning guide further
retrieval~\cite{trivedi2023ircot,asai2024selfrag,jeong2024adaptiverag,yan2025gam}.
They can revise the search as evidence accumulates, but a model's judgment
that the evidence is sufficient does not identify which facts support
the answer or how they connect.

We argue that these relational requirements should be stated and executed
explicitly. Conjunctive queries offer a natural formalism: a query
specifies relations whose arguments must admit a consistent
assignment~\cite{chandra1977conjunctive}, and a \emph{witness} is a set of
facts that realizes one such assignment, thereby explaining the answer
through its provenance~\cite{buneman2001why,green2007provenance}. Applied
to memory, this perspective separates two tasks of different nature. The
linguistic task of deciding \emph{what} to find is left to a language
model; the combinatorial task of finding facts that jointly meet that
specification is delegated to an executor.

Building on this idea, we introduce \sys, a retrieval framework for
long-term agent memory. \sys organizes the interaction history into a
graph of dated propositions that remain linked to their source dialogue.
Given a question, a language model writes a conjunctive plan grounded in
retrieved evidence, and an executor searches the graph for witnesses of
that plan. When search fails or a candidate is rejected, the unsatisfied
requirement tells the planner what to revise. Accepted witnesses finally
guide the selection of the evidence passed to the reader. Because each
answer is tied to an explicit derivation over stored facts, retrieval
becomes inspectable: one can see not only what was retrieved, but why it
supports the answer.

This paper makes the following contributions:
\begin{itemize}
  \item We formulate long-term memory retrieval as \emph{witness search}
    over conjunctive plans, separating the interpretation of a question
    from the execution of its relational requirements.
  \item We present \sys, a retrieval controller that combines
    evidence-grounded planning, reference-relative temporal ranking,
    consistent-binding search, and verification with gap-driven
    replanning.
  \item We show how witnesses can guide the construction of a compact,
    source-linked reader context composed of dated propositions and
    passage summaries.
\end{itemize}
% Add empirical contributions only after validating the completed runs,
% comparison protocol, uncertainty estimates, and reported configuration.

The remainder of the paper is organized as follows.
\Cref{sec:background} reviews the formal background and related work.
\Cref{sec:problem} formulates the task and \cref{sec:method} presents
\sys. \Cref{sec:setup,sec:results} describe the evaluation, and
\cref{sec:discussion,sec:conclusion} discuss limitations and conclude
the paper.
